\section*{Capítulo \ref{ch:b3:conformal} --- Aplicaciones conformes y
el teorema de la aplicación de Riemann}

\begin{solution}{exo:b3:conformal:1}
(a) $\varphi$ y $\psi(w) = \iu\frac{1+w}{1-w}$ se componen dando la
identidad en ambos órdenes (cálculo directo); $\varphi$ lleva
$\mathbb H$ dentro de $\mathbb D$ y $\psi$ de vuelta
(el~\cref{ex:b3:conformal:mobius}). Valores: $\varphi(\iu) = 0$,
$\varphi(0) = -1$, $\varphi(1) = \frac{1 - \iu}{1 + \iu} =
-\iu$ y $\varphi(z) \to 1$ cuando $z \to \infty$.

(b) Las circunferencias y rectas son los conjuntos de soluciones de
$\alpha\abs z^2 + \bar\beta z + \beta\bar z + \gamma = 0$
($\alpha, \gamma \in \R$, $\beta \in \C$, $\abs\beta^2 >
\alpha\gamma$): $\alpha \neq 0$ circunferencias, $\alpha = 0$ rectas.
Sustituyendo $z = 1/w$ y multiplicando por $\abs w^2$:
$\gamma\abs w^2 + \beta w + \bar\beta\bar w + \alpha = 0$ ---
la misma familia. Las aplicaciones afines conservan claramente la
familia, y toda \omterm{ex:b3:conformal:mobius}{transformación de Möbius} es composición de aplicaciones
afines y una inversión ($\frac{az+b}{cz+d} = \frac ac + \frac{bc -
ad}{c}\cdot\frac1{cz + d}$ para $c \neq 0$).
\end{solution}

\begin{solution}{exo:b3:conformal:2}
(a) Schwarz da $\abs{f(a)} \leq \abs a$ con igualdad (ambos miembros
valen $= \abs a$): el caso de igualdad obliga a $f(z) =
\eu^{\iu\theta}z$, y $f(a) = a$ fija $\eu^{\iu\theta} = 1$.

(b) Sean $a \neq b$ los puntos fijos y $g =
\varphi_a\circ f\circ\varphi_{-a} \in
\operatorname{Aut}(\mathbb D)$ --- usando
el~\cref{thm:b3:conformal:autdisc} para $\varphi_{\pm a}$. Entonces
$g(0) = \varphi_a(f(a)) = 0$ y $g(c) = c$ para $c =
\varphi_a(b) \neq 0$: por (a), $g = \mathrm{id}$, de modo que $f =
\varphi_{-a}\circ\varphi_a = \mathrm{id}$.
\end{solution}

\begin{solution}{exo:b3:conformal:3}
Fíjese $z$ y póngase $g = \varphi_{f(z)}\circ f\circ\varphi_{-z}$:
$\mathbb D \to \mathbb D$ \omterm{def:b3:holomorphic:holo}{holomorfa}
con $g(0) = 0$, de modo que $\abs{g'(0)} \leq 1$ (Schwarz). Regla de la
cadena con
$\varphi_a'(\zeta) = \frac{1 - \abs a^2}{(1 - \bar
a\zeta)^2}$:
\[
g'(0) = \varphi_{f(z)}'\bigl(f(z)\bigr)\cdot f'(z)\cdot
\varphi_{-z}'(0)
= \frac{1}{1 - \abs{f(z)}^2}\cdot f'(z)\cdot(1 - \abs z^2),
\]
de donde la desigualdad de Schwarz--Pick. La igualdad en algún $z$ hace
de $g$ una rotación, luego $f =
\varphi_{-f(z)}\circ(\text{rotación})\circ\varphi_z \in
\operatorname{Aut}(\mathbb D)$ --- y entonces la igualdad vale en todas
partes (calcúlese, o reaplíquese intercambiando los papeles de
$f, f^{-1}$). Las autoaplicaciones
\omterm{def:b3:holomorphic:holo}{holomorfas} del disco son
$1$-lipschitzianas para la métrica hiperbólica $\frac{2\abs{\dd
z}}{1 - \abs z^2}$; los automorfismos son sus isometrías.
\end{solution}

\begin{solution}{exo:b3:conformal:4}
(a) $z \mapsto \eu^z$: lleva $\{0 < \operatorname{Im}z <
\pi\}$ biyectivamente sobre $\mathbb H$ ($\eu^{x+\iu y} =
\eu^x\eu^{\iu y}$: \omterm{def:b3:modules:free}{módulo libre}, argumento $y \in
\intoo0\pi$), es \omterm{def:b3:holomorphic:holo}{holomorfa} con
derivada no nula e inversa \omterm{def:b3:holomorphic:holo}{holomorfa}
(el $\log$ principal). (b) $z \mapsto z^2$ duplica los argumentos: el
cuadrante abierto $\{0 < \arg z < \frac\pi2\}$ va
\omterm{def:b3:conformal:conformal}{conformemente} sobre $\mathbb
H$ (inversa: raíz cuadrada principal). (c)
$z \mapsto \frac{1 + z}{1 - z}$ lleva $\mathbb D$ sobre el semiplano
derecho y conserva la simetría superior/inferior: envía el semidisco
superior sobre el primer cuadrante; después elévese al cuadrado, por
(b), para llegar a $\mathbb H$: $z \mapsto \bigl(\frac{1
+ z}{1 - z}\bigr)^2$. (d) El \omterm{thm:b3:conformal:autdisc}{factor de
Blaschke} $\varphi_{1/2}(z) = \frac{z -
\frac12}{1 - \frac z2}$: $\varphi_{1/2}(\tfrac12) = 0$ y
$\varphi_{1/2}'(\tfrac12) = \frac{1 - \frac14}{(1 -
\frac14)^2} = \frac43 > 0$.
\end{solution}

\begin{solution}{exo:b3:conformal:5}
(a) Una $\C \to \mathbb D$ conforme (o $\mathbb H$, tras componer con
Cayley) es una función entera acotada: constante por Liouville --- no
biyectiva. Una $\mathbb
D \to \C$ conforme tendría inversa conforme $\C \to \mathbb D$: la misma
contradicción.

(b) El cuadrado es convexo, luego
\omterm{def:b3:conformal:simplyconnected}{simplemente conexo}, y propio:
\omterm{def:b3:conformal:conformal}{conformemente} $\mathbb D$
(el~\cref{thm:b3:conformal:rmt}). El plano hendido
$\C\setminus\intoc{-\infty}0$ es estrellado respecto de $1$ (los
segmentos desde $1$ evitan el corte), luego
\omterm{def:b3:conformal:simplyconnected}{simplemente conexo}, y propio:
\omterm{def:b3:conformal:conformal}{conformemente} $\mathbb D$. El disco
punteado: si $g \colon \mathbb D\setminus\{0\} \to \mathbb D$ fuera
conforme, $g$ está acotada, de modo que $0$ es evitable
(el~\cref{thm:b3:residues:riemanncw}): $g$ se extiende a $\tilde g
\colon \mathbb D \to \mathbb D$, y $\tilde g(0)$, al estar en la imagen
abierta $g(\mathbb D\setminus\{0\}) = \mathbb D$, es también $g(w)$ para
cierto $w \neq 0$; dos \omterm{def:b3:topology:topology}{entornos}
disjuntos de $0$ y $w$ tienen imágenes abiertas que comparten el valor
$\tilde g(0)$, luego comparten \emph{otros} valores también
(abiertos): $g$ toma algún valor dos
veces en $\mathbb D\setminus\{0\}$ --- en contra de la inyectividad. La
corona $A = \{1 < \abs z < 2\}$: supóngase $F \colon \mathbb D
\to A$ conforme. $F$ no tiene ceros en el
\omterm{def:b3:conformal:simplyconnected}{simplemente conexo}
$\mathbb D$, de modo que $F = \eu^L$ con $L$
\omterm{def:b3:holomorphic:holo}{holomorfa}
(la~\cref{def:b3:conformal:simplyconnected}). Sea $\sigma$ la
circunferencia $\abs z = \frac32$ contenida en $A$ y $\gamma =
F^{-1}\circ\sigma$, un \omterm{def:b3:topology:pathconnected}{camino}
cerrado en $\mathbb D$; entonces
\[
1 = \operatorname{Ind}_\sigma(0)
= \frac1{2\iu\pi}\int_{F\circ\gamma}\frac{\dd w}{w}
= \frac1{2\iu\pi}\int_\gamma\frac{F'}{F}
= \frac1{2\iu\pi}\int_\gamma L' = 0
\]
($L'$ tiene primitiva): contradicción. Ni el disco punteado ni la corona
son un disco disfrazado.
\end{solution}

\begin{solution}{exo:b3:conformal:6}
(a) $T(w) = \frac{w - 1}{w + 1}$ lleva $\{\operatorname{Re}w >
0\}$ \omterm{def:b3:conformal:conformal}{conformemente} sobre
$\mathbb D$ (Cayley rotada: $\abs{w -
1} < \abs{w + 1}$ si y solo si $\operatorname{Re}w > 0$), con $T(1) =
0$. Para $f \in \mathcal F$, $g = T\circ f\colon \mathbb D \to
\mathbb D$ es \omterm{def:b3:holomorphic:holo}{holomorfa} con
$g(0) = 0$: Schwarz da
$\abs{g(z)} \leq \abs z$.

(b) Invirtiendo $T$: $f = \frac{1 + g}{1 - g}$, de modo que
\[
\abs{f(z)} \leq \frac{1 + \abs{g(z)}}{1 - \abs{g(z)}}
\leq \frac{1 + \abs z}{1 - \abs z} :
\]
localmente acotada (uniformemente en $\abs z \leq r < 1$). La igualdad
en $z_0 \neq 0$ obliga a $\abs{g(z_0)} = \abs{z_0}$ y al alineamiento:
$g$ es una rotación, es decir, $f(z) = \frac{1 +
\eu^{\iu\theta}z}{1 - \eu^{\iu\theta}z}$ --- los \emph{extremales de
Herglotz}, aplicaciones conformes
sobre el semiplano derecho.
\end{solution}

\begin{solution}{exo:b3:conformal:7}
(i) La \omterm{def:b3:topology:connected}{conexión} simple interviene exactamente dos veces, a través de la
existencia de raíces cuadradas
\omterm{def:b3:holomorphic:holo}{holomorfas} de funciones sin ceros
(la~\cref{def:b3:conformal:simplyconnected}): en el Paso 0 (la raíz de
$z - b$) y en el Paso 2 (la raíz de $\varphi_a\circ f$). (ii)
$\Omega \neq \C$ proporciona el punto $b$ del Paso 0 --- sin él, la
familia $\mathcal F$ carece de aplicaciones acotadas inyectivas
(Liouville). (iii) La \omterm{def:b3:topology:connected}{conexión} se
usa siempre que hablan el principio de identidad o Hurwitz
(el~\cref{exo:b3:residues:8}): el límite extremal es «inyectivo o
constante», y la constancia queda excluida por $s > 0$; también en
«derivada nula implica constante». Fallo para $\C$: el Paso 0 es
imposible, y la conclusión es falsa (el~\cref{exo:b3:conformal:5}(a)).
Fallo para la corona: no es
simplemente \omterm{def:b3:topology:connected}{conexa} --- la
construcción de la raíz cuadrada se rompe (por ejemplo, la propia $z$,
sin ceros en $A$, no tiene raíz cuadrada
\omterm{def:b3:holomorphic:holo}{holomorfa}: el mismo cálculo de índice
que en el~\cref{exo:b3:conformal:5}(b) con
$\frac12\int_\sigma\frac{\dd z}z \notin 2\iu\pi\Z$) --- y la conclusión
también es falsa.
\end{solution}

\begin{solution}{exo:b3:conformal:8}
Sustituyendo el desarrollo de Fourier de $g$ en la integral de Poisson y
usando $\frac1{2\pi}\int P_r(\varphi - t)\eu^{\iu nt}\dd t =
r^{\abs n}\eu^{\iu n\varphi}$ (léase en la serie de $P_r$):
$u(r\eu^{\iu\varphi}) = \sum_nc_n(g)\,r^{\abs
n}\eu^{\iu n\varphi}$, con el intercambio justificado por la
convergencia normal ($\abs{c_n} \leq \norm g_\infty$, $r < 1$).

(a) $g = \cos t$: $c_{\pm1} = \frac12$, de modo que $u =
r\cos\varphi = \operatorname{Re}z = x$ --- en efecto armónica y con los
valores frontera correctos.

(b) $\cos^2t = \frac12 + \frac{\cos 2t}2$: $u = \frac12 +
\frac{r^2\cos2\varphi}2 = \frac12 +
\frac{\operatorname{Re}(z^2)}2 = \frac12 + \frac{x^2 -
y^2}{2}$.

(c) $g = \mathbf 1_{(0,\pi)}$ (semicircunferencia superior): $c_0 =
\frac12$ y $c_n = \frac{1 - (-1)^n}{2\iu\pi n}$ para $n \neq
0$, de modo que
\[
u(r\eu^{\iu\varphi}) = \frac12 + \frac2\pi\sum_{k\geq0}
\frac{r^{2k+1}\sin\bigl((2k+1)\varphi\bigr)}{2k + 1},
\qquad u(0) = \frac12 :
\]
el \omterm{ex:b3:groups:actions}{centro} ve exactamente el promedio de
los datos frontera --- la propiedad del valor medio en persona.
\end{solution}

\begin{solution}{exo:b3:conformal:9}
A partir de $(1-r)^2 \leq 1 - 2r\cos\theta + r^2 \leq (1+r)^2$:
\[
\frac{1-r}{1+r} = \frac{1 - r^2}{(1+r)^2} \leq P_r(\theta)
\leq \frac{1 - r^2}{(1 - r)^2} = \frac{1+r}{1-r} .
\]
Para $u \geq 0$ armónica en $\mathbb D$ y $s < 1$: $u_s(z) =
u(sz)$ es armónica en un \omterm{def:b3:topology:topology}{entorno} de
$\bar{\mathbb D}$, luego coincide con su integral de Poisson
(el~\cref{thm:b3:conformal:poisson}, unicidad, aplicado a sus propios
valores frontera); encajonando el núcleo y usando el valor medio
$\frac1{2\pi}\int u_s(\eu^{\iu t})\dd t = u(0)$:
\[
\frac{1-r}{1+r}\,u(0) \leq u(s\,r\eu^{\iu\varphi}) \leq
\frac{1+r}{1-r}\,u(0).
\]
Hágase $s \to 1^-$ con $r\eu^{\iu\varphi}$ fijo
(\omterm{def:b3:topology:continuity}{continuidad} de $u$): las
desigualdades de Harnack. Si $u$ es armónica en $\C$ con $u \geq m$:
aplíquese Harnack a $u - m$ en discos $D(0, R)$, es decir, a
$z \mapsto u(Rz) - m$: para $z$ fijo y $r =
\abs z/R \to 0$, ambas cotas tienden a $u(0) - m$: $u(z) =
u(0)$ --- constante (un Liouville bilátero a partir de una cota
unilátera).
\end{solution}

\begin{solution}{exo:b3:conformal:10}
Localmente, $u = \operatorname{Re}F$ con $F$
\omterm{def:b3:holomorphic:holo}{holomorfa}
(la~\cref{prop:b3:conformal:harmonicholo}), de modo que $u\circ f =
\operatorname{Re}(F\circ f)$ es armónica donde esté definida: la
armonicidad es \omterm{def:b3:conformal:conformal}{conformemente}
invariante. En $\mathbb H$: el logaritmo principal da
$\log z = \ln\abs z + \iu\arg z$
\omterm{def:b3:holomorphic:holo}{holomorfa} en $\mathbb H$, de modo que
$u = \frac1\pi\arg z =
\operatorname{Im}\bigl(\frac1\pi\log z\bigr)$ es armónica, con límites
frontera: para $x > 0$, $\arg \to 0$, $u \to 0$; para $x < 0$,
$\arg \to \pi$, $u \to 1$: los datos $\mathbf
1_{\intoo{-\infty}0}$ en todo $x \neq 0$. Transportándola por la
\omterm{ex:b3:conformal:mobius}{transformación de Cayley} (que envía
$\mathbb D \to \mathbb H$ tras una inversión y hace corresponder la
semicircunferencia superior con el semieje negativo, salvo la rotación
que se fija siguiendo tres puntos frontera), $u\circ(\text{Cayley})$
resuelve el problema en el disco del~\cref{exo:b3:conformal:8}(c);
evaluando en el \omterm{ex:b3:groups:actions}{centro} se recupera allí
$u = \frac12$, y la forma cerrada $\frac1\pi\arg$ puede contrastarse con
la serie sumando identidades de tipo
$\sum\frac{r^{2k+1}\sin((2k+1)\varphi)}{2k+1} =
\frac12\arctan\frac{2r\sin\varphi}{1 - r^2}$ --- la vía elemental hacia
la misma respuesta.
\end{solution}

\begin{solution}{exo:b3:conformal:11}
(a) Sea $a \neq b$ fijo. Conjugando por $\varphi_a(z) =
\frac{z - a}{1 - \bar az}$ (un automorfismo que intercambia $a$ y $0$),
$g = \varphi_a\circ f\circ\varphi_a^{-1}$ fija $0$ y el punto
$c = \varphi_a(b) \neq 0$. Schwarz: $\abs{g(z)} \leq \abs z$, y en
$z = c$ se da la igualdad ($g(c) = c$): el caso de igualdad obliga a
$g(z) = \lambda z$ con $\abs\lambda = 1$, y $\lambda c = c$ da
$\lambda =
1$: $g = \mathrm{id}$, luego $f = \mathrm{id}$.

(b) $h(z) = f(z)/z$ (\omterm{thm:b3:residues:riemanncw}{singularidad
evitable} en $0$) es \omterm{def:b3:holomorphic:holo}{holomorfa} en
$\mathbb D$ con $\abs h \leq 1$ (Schwarz); $\abs h < 1$ en todas partes,
pues en caso contrario el principio del máximo (máximo
\omterm{def:b3:topology:interior}{interior} de $\abs h$) haría de $h$
una constante unimodular, es decir, de $f$ una rotación --- excluido. En
el \omterm{def:b3:topology:compact}{compacto} $\bar D(0,r)$,
$c_r = \max\abs h < 1$, de modo que allí $\abs{f(z)} \leq
c_r\abs z$; además, $f$ lleva $\bar D(0,r)$ dentro de sí mismo
($c_r\abs z \leq r$), de modo que la cota itera:
$\abs{f^{\circ n}(z)} \leq c_r^n\,r \to 0$ uniformemente en
$\bar D(0, r)$.

(c) Puntos fijos de $\frac{z^2 + z}2$: $z^2 + z = 2z$ si y solo si
$z(z - 1) = 0$; solo $z = 0$ está en $\mathbb D$ ($z = 1$ está en la
frontera). No es una rotación ($f'(0) = \frac12$), de modo que las
\omterm{def:b3:groups:action}{órbitas} tienden a $0$; cuantitativamente,
$f(z) = \frac z2(1 + z)$ da $\abs{f(z)} \leq \frac{3}{4}\abs z$ en
$\abs z \leq
\frac12$ y, una vez que la \omterm{def:b3:groups:action}{órbita} es
pequeña, $\abs{f(z)} \approx
\frac{\abs z}2$: asintóticamente geométrica de razón $f'(0) = \frac12$.
A partir de $z_0 = \frac12$: $z_1 = \frac38$, $z_2 \approx 0.258$,
$z_3 \approx 0.162$ --- se reduce a la mitad en cada paso, como predice
el multiplicador.
\end{solution}

\begin{solution}{exo:b3:conformal:12}
(a) $\Delta u = 6x - 6x + 0 = 0$. Cauchy--Riemann exige
$v_y = u_x = 3x^2 - 3y^2$ y $v_x = -u_y = 6xy - 2$. Integrando la
primera en $y$: $v = 3x^2y - y^3 + c(x)$; sustituyendo en la segunda:
$6xy + c'(x) = 6xy - 2$, de modo que $c(x)
= -2x + C$. Así, $v = 3x^2y - y^3 - 2x + C$ y
\[
f = u + \iu v = (x^3 - 3xy^2) + \iu(3x^2y - y^3)
+ 2y - 2\iu x + \iu C = z^3 - 2\iu z + \iu C .
\]

(b) La forma $\omega = -u_y\,\dd x + u_x\,\dd y$ es cerrada precisamente
porque $\Delta u = 0$ ($\partial_y(-u_y) =
-u_{yy} = u_{xx} = \partial_x(u_x)$). En un
abierto estrellado, el lema de
Poincaré (el~\cref{thm:b3:forms:poincare}; o la construcción de
primitivas del~\cref{thm:b3:holomorphic:cauchy} aplicada a la
\omterm{def:b3:holomorphic:holo}{holomorfa} $u_x - \iu u_y$) proporciona
$v$ con $\dd v = \omega$, es decir, el sistema de Cauchy--Riemann:
$u + \iu v$ es \omterm{def:b3:holomorphic:holo}{holomorfa}. Dos
conjugadas difieren en una función de gradiente nulo: una constante
(\omterm{def:b3:topology:connected}{conexión}).

(c) Para $u = \ln\abs z$ en $\C^*$: $\omega = -u_y\dd x +
u_x\dd y = \frac{-y\,\dd x + x\,\dd y}{x^2 + y^2} =
\omega_\theta$, la forma angular (el~\cref{ex:b3:forms:angular}), cuya
integral a lo largo de la circunferencia unidad vale $2\pi \neq 0$: no
es exacta, de modo que no existe conjugada global --- una conjugada
sería una determinación \omterm{def:b3:topology:continuity}{continua}
del argumento, y el índice es exactamente la obstrucción. Localmente (en
cualquier subdominio estrellado) sirve $v =
\arg z$ y $u + \iu v = \log z$: el fallo es global, no local.
\end{solution}

\begin{solution}{pb:b3:conformal:1}
\textbf{1.} $g$ es inyectiva y
\omterm{def:b3:holomorphic:holo}{holomorfa}; en $C_\rho$ ($\rho > 1$) es
regular, y el área encerrada es
$\frac1{2\iu}\oint_{g(C_\rho)}\bar\zeta\,\dd\zeta$ (la fórmula del área
de Green--Riemann de segundo año, aplicada con orientación positiva).
Sustituyendo $\zeta = g(w)$, $w =
\rho\eu^{\iu\theta}$:
\[
A_\rho = \frac1{2\iu}\int_{C_\rho}\overline{g(w)}\,g'(w)\dd
w = \frac\rho2\int_0^{2\pi}\overline{g(\rho\eu^{\iu\theta})}
\,g'(\rho\eu^{\iu\theta})\,\eu^{\iu\theta}\,\dd\theta .
\]
Insértense $\bar g = \rho\eu^{-\iu\theta} + \bar b_0 +
\sum_m\bar b_m\rho^{-m}\eu^{\iu m\theta}$ y $g' = 1 -
\sum_nnb_n\rho^{-n-1}\eu^{-\iu(n+1)\theta}$: tras multiplicar por
$\eu^{\iu\theta}$, solo los productos de frecuencia cero sobreviven a la
integración en $\theta$ (la convergencia normal justifica el trabajo
término a término): el par $(\rho\eu^{-\iu\theta})\cdot1$ aporta
$2\pi\rho$, y cada par $\bar b_n\rho^{-n}\eu^{\iu
n\theta}\cdot(-nb_n\rho^{-n-1}\eu^{-\iu(n+1)\theta})$ aporta
$-2\pi n\abs{b_n}^2\rho^{-2n-1}$:
\[
A_\rho = \pi\rho^2 - \pi\sum_{n\geq1}n\abs{b_n}^2\rho^{-2n}.
\]

\textbf{2.} Las áreas son no negativas: $\sum_{n\leq
N}n\abs{b_n}^2\rho^{-2n} \leq \rho^2$ para todo $N$; hágase
$\rho \to 1^+$ y después $N \to \infty$:
$\sum_{n\geq1}n\abs{b_n}^2 \leq 1$. La igualdad en $\abs{b_1}
\leq 1$ obliga a que todos los demás $b_n = 0$: $g(w) = w + b_0 +
\eu^{\iu\alpha}/w$, que aplica sobre el complementario de un segmento de
longitud $4$ (una aplicación de tipo Joukowski): los extremales.

\textbf{3.} $f(z)/z = 1 + a_2z + \cdots$ es
\omterm{def:b3:holomorphic:holo}{holomorfa} y sin ceros en $\mathbb D$
($f$ solo se anula en $0$, y simplemente: inyectividad), luego también
lo es $f(z^2)/z^2$, que tiene una raíz cuadrada
\omterm{def:b3:holomorphic:holo}{holomorfa} $\varphi$ con
$\varphi(0) = 1$ (la~\cref{def:b3:conformal:simplyconnected};
$\mathbb D$ es convexo). Entonces $h(z) = z\varphi(z^2)$ cumple $h^2 =
f(z^2)$, $h(z) = z\bigl(1 + \frac{a_2}2z^2 + \cdots\bigr)$ (serie
binómica de la raíz), y $h(z) = z\varphi(z^2)$ es impar por construcción
($\varphi(z^2)$ es par). Inyectividad: $h(z) = h(z')$ da $f(z^2) =
f(z'^2)$, de modo que $z^2 = z'^2$, es decir, $z' = \pm z$; si
$z' = -z$, la imparidad da $h(z) = -h(z)$, de modo que $h(z) = 0$, lo
que fuerza $z
= 0$ ($\varphi$ sin ceros): $z = z' = 0$.

\textbf{4.} $g(w) = 1/h(1/w)$: para $\abs w > 1$, $1/w \in
\mathbb D\setminus\{0\}$ y allí $h \neq 0$: bien definida, inyectiva
(composición de inyecciones), con desarrollo
\[
g(w) = \frac{1}{\frac1w\bigl(1 + \frac{a_2}{2w^2} +
\cdots\bigr)} = w\Bigl(1 - \frac{a_2}{2w^2} + \cdots\Bigr)
= w - \frac{a_2}{2}\,w^{-1} + \cdots :
\]
de la forma de la Parte I con $b_1 = -\frac{a_2}2$. El teorema del área
da $\abs{\frac{a_2}2} \leq 1$: $\abs{a_2} \leq 2$.

\textbf{5.} $k(z) = \frac z{(1-z)^2} = z\sum_{m\geq0}(m +
1)z^m = \sum_{n\geq1}nz^n$: coeficientes $a_n = n$, de modo que $a_2 =
2$. Univalencia e imagen: $k = \frac14\bigl[w^2 - 1\bigr]$ con
$w = \frac{1 + z}{1 - z}$, una
\omterm{def:b3:conformal:conformal}{aplicación conforme} de $\mathbb D$
sobre el semiplano derecho; $w^2$ lleva ese semiplano
\omterm{def:b3:conformal:conformal}{conformemente} sobre
$\C\setminus\intoc{-\infty}0$; después $\frac{\cdot - 1}4$ da
$\C\setminus\intoc{-\infty}{-\frac14}$: inyectiva en cada etapa, con la
imagen anunciada.

\textbf{6.} $F = \frac{cf}{c - f}$: como $c \notin f(\mathbb
D)$, el denominador no se anula nunca: $F$ es
\omterm{def:b3:holomorphic:holo}{holomorfa} e inyectiva
($w \mapsto \frac{cw}{c - w}$ es de Möbius, inyectiva fuera de $w = c$).
Desarrollo: con $f = z + a_2z^2 + \cdots$,
\[
F = f\cdot\frac1{1 - f/c} = \bigl(z + a_2z^2\bigr)\Bigl(1 +
\frac zc\Bigr) + O(z^3)
= z + \Bigl(a_2 + \frac1c\Bigr)z^2 + O(z^3):
\]
$F \in \mathcal S$ con $A_2 = a_2 + \frac1c$.

\textbf{7.} Bieberbach dos veces: $\abs{a_2} \leq 2$ y
$\abs{a_2 + \frac1c} \leq 2$, de modo que $\abs{\frac1c} \leq 4$:
$\abs c \geq \frac14$. Todo valor omitido queda fuera de
$D(0,\frac14)$, i.e.\ $f(\mathbb D) \supseteq D(0, \frac14)$.
Óptimo: la función de Koebe omite $-\frac14$ (pregunta 5).

\textbf{8.} Sean $d = d(z_0, \partial\Omega)$ y $\psi =
\varphi^{-1} \colon \mathbb D \to \Omega$, $\psi(0) = z_0$,
$\psi'(0) = 1/\varphi'(z_0) > 0$. La normalización $\tilde
f(z) = \frac{\psi(z) - z_0}{\psi'(0)}$ está en $\mathcal S$, de modo que
su imagen contiene $D(0, \frac14)$; deshaciendo la escala,
$\Omega = \psi(\mathbb D) \supseteq D\bigl(z_0,
\tfrac{\psi'(0)}4\bigr)$, luego $d \geq \frac{\psi'(0)}4 =
\frac1{4\varphi'(z_0)}$: la desigualdad izquierda. Para la derecha:
$\varphi\circ(z_0 + d\,\cdot)$ lleva $\mathbb D$ dentro de $\mathbb D$
con $0 \mapsto 0$; Schwarz acota su derivada en $0$:
$d\,\varphi'(z_0) \leq 1$, es decir, $\frac1{\varphi'(z_0)} \geq d$ ---
y juntas
\[
d \leq \frac1{\varphi'(z_0)} \leq 4d .
\]
(La desigualdad izquierda tal como aparece en el enunciado es la misma
cadena reordenada.)

\textbf{9.} Para la función de Koebe $a_n = n$: igualdad en toda la
conjetura; sus rotaciones $\eu^{-\iu\theta}k(\eu^{\iu\theta}z) = \sum
n\eu^{\iu(n-1)\theta}z^n$ tienen también $\abs{a_n} = n$. Llevando la
pregunta 4 más lejos: $h = z + \frac{a_2}2z^3 + c_5z^5 + \cdots$ con
$h^2 = f(z^2)$; comparando los coeficientes de $z^6$: $2c_5 +
\frac{a_2^2}4 = a_3$, de modo que $c_5 = \frac{a_3}2 -
\frac{a_2^2}8$. Entonces
\[
g(w) = \frac1{h(1/w)} = w\Bigl(1 - \frac{a_2}{2w^2} +
\Bigl(\frac{a_2^2}4 - c_5\Bigr)\frac1{w^4} + \cdots\Bigr)
= w - \frac{a_2}2w^{-1} + \Bigl(\frac{3a_2^2}8 -
\frac{a_3}2\Bigr)w^{-3} + \cdots,
\]
y el teorema del área ($\sum n\abs{b_n}^2 \leq 1$) produce
\[
\Bigl|\frac{a_2}2\Bigr|^2 + 3\,\Bigl|\frac{3a_2^2}8 -
\frac{a_3}2\Bigr|^2 \leq 1 .
\]
Comprobación con Koebe ($a_2 = 2$, $a_3 = 3$): $\abs{\frac{a_2}2}^2 = 1$
y $\frac{3\cdot4}8 - \frac32 = 0$: total exactamente $1$ --- extremal,
como debía ser.

\textbf{10.} $\varphi$ es un automorfismo del disco con
$\varphi(0) = z_0$, de modo que $f\circ\varphi$ es univalente en
$\mathbb D$ (composición de inyecciones), y $f'(z_0) \neq
0$ (la~\cref{def:b3:conformal:conformal}): $F$ está bien definida y es
univalente. $F(0) = 0$. Regla del cociente:
\[
\varphi'(z) = \frac{(1 + \bar z_0z) - (z + z_0)\bar z_0}
{(1 + \bar z_0z)^2}
= \frac{1 - \abs{z_0}^2}{(1 + \bar z_0z)^2},
\qquad \varphi'(0) = 1 - \abs{z_0}^2 ,
\]
de modo que $(f\circ\varphi)'(0) = f'(z_0)(1 - \abs{z_0}^2)$ y
$F'(0) = 1$: $F \in \mathcal S$.

\textbf{11.} Con $G = f\circ\varphi$: $G'' =
f''(\varphi)\,\varphi'^2 + f'(\varphi)\,\varphi''$, y
$\varphi''(z) = -2\bar z_0(1 - \abs{z_0}^2)(1 + \bar
z_0z)^{-3}$ da $\varphi''(0) = -2\bar z_0(1 -
\abs{z_0}^2)$. Por tanto,
\[
A_2 = \frac{G''(0)}{2f'(z_0)(1 - \abs{z_0}^2)}
= \frac12\Bigl[\bigl(1 - \abs{z_0}^2\bigr)
\frac{f''(z_0)}{f'(z_0)} - 2\bar z_0\Bigr] .
\]
Bieberbach para $F$ (pregunta 4): $\abs{A_2} \leq 2$, es decir,
$\bigl|(1 - \abs{z_0}^2)\frac{f''(z_0)}{f'(z_0)} - 2\bar
z_0\bigr| \leq 4$. Multiplíquese todo por $z_0/(1 - \abs{z_0}^2)$, cuyo
módulo es $r/(1 - r^2)$ para $z_0 = r\eu^{\iu\theta}$, y úsese $z_0\bar
z_0 = r^2$:
\[
\Bigl|\,z_0\frac{f''(z_0)}{f'(z_0)} -
\frac{2r^2}{1 - r^2}\Bigr| \leq \frac{4r}{1 - r^2} .
\]

\textbf{12.} Un número complejo a distancia $\rho$ del punto real $c$
tiene parte real en $\intcc{c - \rho}{c +
\rho}$: aplíquese esto a $c = \frac{2r^2}{1-r^2}$, $\rho =
\frac{4r}{1-r^2}$.

\textbf{13.} Para $\mathcal C^1$ $g$ no nula:
$\log\abs g = \frac12\log(g\bar g)$, de modo que $\frac{\dd}{\dd
t}\log\abs g = \frac{g'\bar g + g\bar g'}{2\abs g^2} =
\operatorname{Re}\frac{g'}g$. Con $g(t) =
f'(t\eu^{\iu\theta})$ (sin ceros: univalencia), $g'(t) =
\eu^{\iu\theta}f''(t\eu^{\iu\theta})$, de modo que en $z =
t\eu^{\iu\theta}$:
\[
\frac{\dd}{\dd t}\log\bigl|f'(t\eu^{\iu\theta})\bigr|
= \operatorname{Re}\Bigl(\eu^{\iu\theta}
\frac{f''}{f'}\Bigr)
= \frac1t\operatorname{Re}\Bigl(z\frac{f''}{f'}\Bigr)
\in \Bigl[\frac{2t - 4}{1 - t^2},\
\frac{2t + 4}{1 - t^2}\Bigr]
\]
por la pregunta 12 con radio $t$. Como $\frac{\dd}{\dd
t}\log\frac{1+t}{(1-t)^3} = \frac1{1+t} + \frac3{1-t} =
\frac{2t+4}{1-t^2}$ y $\frac{\dd}{\dd
t}\log\frac{1-t}{(1+t)^3} = \frac{-1}{1-t} - \frac3{1+t} =
\frac{2t-4}{1-t^2}$, integrando de $0$ a $r$ (las tres funciones se
anulan en $t = 0$, $f'(0) = 1$) se obtiene
\[
\log\frac{1-r}{(1+r)^3} \leq \log\abs{f'(z)} \leq
\log\frac{1+r}{(1-r)^3} :
\]
el teorema de distorsión, tras exponenciar.

\textbf{14.} Superior: a lo largo del segmento $[0, z]$,
\[
\abs{f(z)} = \Bigl|\int_0^rf'(t\eu^{\iu\theta})
\eu^{\iu\theta}\dd t\Bigr|
\leq \int_0^r\frac{1+t}{(1-t)^3}\dd t
= \frac{r}{(1-r)^2}
\]
($\frac{\dd}{\dd t}\frac t{(1-t)^2} = \frac{1+t}{(1-t)^3}$).
Inferior: obsérvese que $\frac r{(1+r)^2} < \frac14$ para $r < 1$ (dice
$(1-r)^2 > 0$), de modo que si $\abs{f(z)} \geq \frac14$ no hay nada que
demostrar. En caso contrario, el segmento $[0, f(z)]$ está contenido en
$D(0, \frac14) \subseteq f(\mathbb D)$ (pregunta 7), y
$\gamma = f^{-1}\circ[0, f(z)]$ es un
\omterm{def:b3:topology:pathconnected}{camino} $\mathcal C^1$ de $0$ a
$z$ en $\mathbb D$ ($f^{-1}$
\omterm{def:b3:holomorphic:holo}{holomorfa},
el~\cref{cor:b3:residues:openmapping}). Sustituyendo $w =
f(\zeta)$,
\[
\abs{f(z)} = \int_{[0,f(z)]}\abs{\dd w}
= \int_\gamma\abs{f'(\zeta)}\,\abs{\dd\zeta}
\geq \int_0^1\frac{1 - \rho(s)}{(1 + \rho(s))^3}
\,\abs{\gamma'(s)}\,\dd s
\]
con $\rho = \abs\gamma$, usando la cota inferior de distorsión con radio
$\rho(s)$. Como $\abs{\gamma'} \geq \rho'$ (donde esté definida; $\rho$
es lipschitziana) y el factor del integrando es positivo,
\[
\abs{f(z)} \geq \int_0^1
\frac{1 - \rho(s)}{(1 + \rho(s))^3}\,\rho'(s)\,\dd s
= \int_0^{r}\frac{1 - u}{(1 + u)^3}\,\dd u
= \frac{r}{(1 + r)^2}
\]
($\rho(0) = 0$, $\rho(1) = r$; la sustitución solo usa una primitiva,
$\frac{\dd}{\dd u}\frac u{(1+u)^2} =
\frac{1-u}{(1+u)^3}$, no la monotonía).

\textbf{15.} $k'(z) = \frac{\dd}{\dd z}\,z(1 - z)^{-2} =
\frac{1 + z}{(1 - z)^3}$. En $z = r$: $k'(r) =
\frac{1+r}{(1-r)^3}$ y $k(r) = \frac r{(1-r)^2}$ --- ambas cotas
superiores alcanzadas. En $z = -r$: $k'(-r) =
\frac{1-r}{(1+r)^3}$ y $\abs{k(-r)} = \frac r{(1+r)^2}$ --- ambas cotas
inferiores alcanzadas. La función de Koebe estira su semieje positivo al
máximo hacia la frontera lejana y comprime al máximo el radio antipodal
hacia la punta de su corte.

\textbf{16.} $\abs{a_2} = 2$ significa $\abs{b_1} = 1$ para $g(w)
= 1/h(1/w) = w - \frac{a_2}2w^{-1} + \cdots$ (pregunta 4); el teorema
del área $\sum n\abs{b_n}^2 \leq 1$ mata entonces todos los demás
coeficientes: $g(w) = w + b_0 + \eu^{\iu\alpha}/w$ con
$\eu^{\iu\alpha} = -\frac{a_2}2$. Como $h$ es impar, $g(-w) =
1/h(-1/w) = -1/h(1/w) = -g(w)$: $g$ es impar, de modo que $b_0 = 0$.
Invirtiendo, $h(z) = 1/g(1/z) = \frac{z}{1 +
\eu^{\iu\alpha}z^2}$, y $f(z^2) = h(z)^2$ da
\[
f(w) = \frac{w}{(1 + \eu^{\iu\alpha}w)^2}
= \frac{w}{(1 - \eu^{\iu\theta}w)^2}
= \eu^{-\iu\theta}k\bigl(\eu^{\iu\theta}w\bigr),
\qquad \eu^{\iu\theta} = -\eu^{\iu\alpha} .
\]
Recíprocamente, cada rotación tiene $a_2 = 2\eu^{\iu\theta}$ de módulo
$2$: los extremales de Bieberbach son exactamente las funciones de Koebe
rotadas.

\textbf{17.} Si $c$ se omite con $\abs c = \frac14$: la pregunta 6 da
$F \in \mathcal S$ con $A_2 = a_2 +
\frac1c$, de modo que $\frac1c = A_2 - a_2$ y
\[
4 = \Bigl|\frac1c\Bigr| = \abs{A_2 - a_2}
\leq \abs{A_2} + \abs{a_2} \leq 2 + 2 = 4 :
\]
con igualdad en toda la cadena, en particular $\abs{a_2} = 2$. Por la
pregunta 16, $f = \eu^{-\iu\theta}k(\eu^{\iu\theta}\cdot)$, cuyo
conjunto omitido es $\eu^{-\iu\theta}
\intoc{-\infty}{-\frac14}$: el único valor omitido de módulo $\frac14$
es $c = -\eu^{-\iu\theta}/4$. (Comprobación: $a_2 = 2\eu^{\iu\theta}$ y
$\frac1c = -4\eu^{\iu\theta}$, de modo que
$A_2 = -2\eu^{\iu\theta} = -a_2$: la desigualdad triangular queda
saturada por antialineamiento, como debía.)

\textbf{18.} Las \omterm{thm:b3:holomorphic:analytic}{estimaciones de
Cauchy} en la circunferencia $\abs z =
r$ (el~\cref{thm:b3:holomorphic:analytic}), combinadas con el teorema de
crecimiento, dan \[
\abs{a_n} \leq r^{-n}\sup_{\abs z = r}\abs f
\leq r^{1-n}\,(1 - r)^{-2} .
\] Tómese $r = 1 - \frac1n$ ($n \geq
2$): $r^{1-n} = \bigl(1 + \frac1{n-1}\bigr)^{n-1} < \eu$ (sucesión
creciente de límite $\eu$) y $(1 - r)^{-2} =
n^2$: $\abs{a_n} < \eu\,n^2$.

\textbf{19.} $U = f(D(0, r))$ es abierto
(el~\cref{cor:b3:residues:openmapping}) y contiene $0$. Frontera: si
$p \in \partial U$, escríbase $p = \lim f(z_k)$ con $z_k \in D(0, r)$;
una subsucesión da $z_k \to z_\infty
\in \bar D(0, r)$, y $f(z_\infty) = p \notin U$ obliga a
$z_\infty \in \partial D(0, r)$: $\partial U \subseteq
f(\partial D(0, r))$, de modo que todo punto frontera de $U$ tiene
módulo $\geq \frac r{(1+r)^2}$ (pregunta 14). Sean ahora
$\abs w < \frac r{(1+r)^2}$ y supóngase $w \notin U$. El segmento
$[0, w]$ es \omterm{def:b3:topology:connected}{conexo}, corta $U$ (en
$0$) y su complementario (en $w$), de modo que corta $\partial U$; pero
todos sus puntos tienen módulo $\leq \abs w < \frac r{(1+r)^2}$:
contradicción. Luego $D\bigl(0, \frac r{(1+r)^2}\bigr)
\subseteq U$. Optimalidad: $k(-r) = -\frac r{(1+r)^2}$ es la imagen de
un punto frontera de $D(0,r)$, y $k$ es inyectiva, de modo que
$k(-r) \notin k(D(0, r))$: el radio no puede aumentarse. Cuando
$r \to 1^-$, $\frac r{(1+r)^2} \to \frac14$: el teorema del cuarto para
el disco \omterm{def:b3:complete:complete}{completo}.

\textbf{20.} Primero, $\partial f(\mathbb D) \neq \emptyset$: en caso
contrario $f(\mathbb D)$, abierto, cerrado y no vacío, sería todo $\C$,
y $f^{-1} \colon \C \to \mathbb D$ sería una función entera acotada no
constante, en contra del~\cref{cor:b3:holomorphic:liouville}.
Desigualdad izquierda: la transformada de Koebe $F$ de la pregunta 10
está en $\mathcal S$ y $F(\mathbb D) = \frac{f(\mathbb D) -
f(z_0)}{f'(z_0)(1-\abs{z_0}^2)}$; por la pregunta 7 contiene
$D(0, \frac14)$, de modo que
\[
f(\mathbb D) \supseteq f(z_0) + D\Bigl(0,\
\tfrac14\abs{f'(z_0)}\bigl(1 - \abs{z_0}^2\bigr)\Bigr) ,
\]
y todo punto de $\partial f(\mathbb D)$ --- disjunto del
abierto $f(\mathbb D)$ --- está a
distancia $\geq
\frac14(1-\abs{z_0}^2)\abs{f'(z_0)}$ de $f(z_0)$. Desigualdad derecha:
sea $d = d(f(z_0), \partial f(\mathbb D)) <
\infty$. El disco $D(f(z_0), d)$ está contenido en $f(\mathbb D)$: un
segmento de $f(z_0)$ a cualquiera de sus puntos permanece a distancia
$< d$ de $f(z_0)$, de modo que nunca corta $\partial f(\mathbb D)$, y el
argumento de \omterm{def:b3:topology:connected}{conexión} de la pregunta
19 lo mantiene en $f(\mathbb D)$. Entonces
$\hat g(w) = f^{-1}(f(z_0) + dw)$ lleva $\mathbb D$ dentro de
$\mathbb D$ con $\hat g(0) = z_0$, y $\chi = \psi^{-1}\circ\hat g$, con
$\psi(z) = \frac{z +
z_0}{1 + \bar z_0z}$, fija $0$: Schwarz
(el~\cref{thm:b3:conformal:schwarz}) da $\abs{\chi'(0)}
\leq 1$. Como $\chi'(0) = \frac{\hat g'(0)}{\psi'(0)} =
\frac{d}{f'(z_0)\,(1 - \abs{z_0}^2)}$, esto es $d \leq (1 -
\abs{z_0}^2)\abs{f'(z_0)}$.

\textbf{21.} $k(\mathbb D) = \C\setminus
\intoc{-\infty}{-\frac14}$, de modo que $\partial k(\mathbb D) =
\intoc{-\infty}{-\frac14}$, y para el punto real positivo
$k(r) = \frac r{(1-r)^2}$ el punto frontera más próximo es
$-\frac14$:
\[
d = k(r) + \frac14 = \frac{4r + (1-r)^2}{4(1-r)^2}
= \frac{(1+r)^2}{4(1-r)^2} .
\]
Miembro izquierdo de la pregunta 20: $\frac14(1 - r^2)k'(r) =
\frac14\,\frac{(1-r^2)(1+r)}{(1-r)^3} =
\frac{(1+r)^2}{4(1-r)^2} = d$: igualdad. (El miembro derecho vale $4d$:
el factor \omterm{def:b3:complete:complete}{completo} $4$ separa ambos lados, y la función de Koebe se
sitúa exactamente en el fondo.)

\textbf{22.} El principio único es
\omterm{thm:b3:hilbert:parseval}{Parseval}: la univalencia prohíbe los
solapamientos, de modo que el área del complementario de la imagen de
$\{\abs w > \rho\}$, desarrollada en modos de Fourier sobre
circunferencias, es no negativa --- el teorema del área es una identidad
$L^2$ con un signo. Todo lo demás es esa desigualdad transportada: una
raíz cuadrada (pregunta 3) la convierte en $\abs{a_2} \leq 2$; una
reflexión de Möbius respecto de un valor omitido (pregunta 6) convierte
$\abs{a_2} \leq 2$ en el teorema del cuarto; los automorfismos del disco
(pregunta 10) extienden $\abs{a_2}
\leq 2$ a todo el disco como teorema de distorsión; la integración
radial convierte la distorsión en crecimiento, y el crecimiento en
recubrimiento. En cada etapa se propaga también el caso de igualdad, y
siempre aterriza en las funciones de Koebe rotadas --- una sola familia
extremal para toda la teoría, igual que las rotaciones son los únicos
extremales del lema de Schwarz. Una desigualdad
\omterm{def:b3:topology:interior}{interior}, más la rigidez de su caso
de igualdad, gobierna toda la geometría: la \omterm{def:b3:representations:rep}{representación} conforme es
el arte de explotar tales desigualdades.

\textbf{23.} El teorema de crecimiento acota $\abs f \leq
\frac{r}{(1-r)^2}$ uniformemente en cada $\bar D(0, r)$, $r <
1$, para todas las $f \in \mathcal S$ a la vez: localmente acotada, de
modo que $\mathcal S$ es una familia normal
(el~\cref{thm:b3:conformal:montel}). Si $f_n \in \mathcal S \to
f$ localmente uniformemente: $f$ es
\omterm{def:b3:holomorphic:holo}{holomorfa} con $f_n' \to f'$ localmente
uniformemente (Weierstrass), de modo que $f(0) = 0$, $f'(0) = 1$ --- en
particular, $f$ no es constante --- y Hurwitz
(el~\cref{exo:b3:residues:8}) hace inyectivo el límite de aplicaciones
inyectivas: $f \in \mathcal S$. Un funcional
\omterm{def:b3:topology:continuity}{continuo} (como
$f \mapsto \abs{a_2} = \frac{\abs{f''(0)}}2$) sobre una clase
compacta \emph{alcanza} su supremo:
las funciones extremales existen antes de saber cuáles son --- el punto
de partida de todo ataque variacional a los problemas de coeficientes,
el de Bieberbach incluido.

\textbf{24.} Escríbase $g(w) = w + A_2w^2 + O(w^3)$ y compóngase:
\[
z = g(f(z)) = f(z) + A_2f(z)^2 + O(z^3)
= z + (a_2 + A_2)z^2 + O(z^3),
\]
de modo que $A_2 = -a_2$ y $\abs{A_2} \leq 2$ (Bieberbach), con igualdad
si y solo si $\abs{a_2} = 2$, es decir, si y solo si $f$ es una función
de Koebe rotada (pregunta 16) --- y entonces $g$ es la inversa
correspondiente, definida en el plano hendido.

\textbf{25.} $f(z_1) - f(z_2) = (z_1 - z_2)\bigl(1 + a(z_1 +
z_2)\bigr)$. Si $\abs a \leq \frac12$: para $z_1 \neq z_2$ en
$\mathbb D$, $\abs{a(z_1 + z_2)} < 2\abs a \leq 1$ (estricto:
$\abs{z_1 + z_2} < 2$), de modo que el segundo factor no puede anularse:
inyectiva, y $f \in \mathcal S$ (las normalizaciones están
incorporadas). Si $\abs a > \frac12$: el punto $s =
-\frac1a$ cumple $\abs s < 2$, de modo que $z_{1,2} = \frac s2 \pm
\varepsilon$ están en $\mathbb D$ para $\varepsilon > 0$ pequeño, son
distintos, y $z_1 + z_2 = s$ mata el factor: $f(z_1)
= f(z_2)$, no inyectiva. Así, $\mathcal S$ contiene $z +
az^2$ exactamente para $\abs a \leq \frac12$. La cota de Bieberbach
$\abs{a_2} \leq 2$ queda, pues, salvajemente sin saturar por los
polinomios de grado $2$ --- los coeficientes $a_n = n$ de la función de
Koebe provienen de una serie infinita que conspira a lo largo del radio
omitido, comportamiento que ningún polinomio (que solo pertenece a
$\mathcal S$ con coeficientes minúsculos) puede imitar.
\end{solution}
